\documentclass[a4paper,14pt]{article}

% preambule.tex

% Русский язык через polyglossia
\usepackage{polyglossia}
\setmainlanguage{russian}
\newfontfamily{\cyrillicfonttt}{Courier New}

% Шрифты
\usepackage{fontspec}
\setmainfont{Times New Roman}

% Графика
\usepackage{graphicx}

% Фиксация float окружений
\usepackage{float}

% Математика
\usepackage{amsmath,amssymb,amsfonts,amsthm}

% Таблицы
\usepackage{makecell}
\usepackage{multirow}

% Листинги кода
\usepackage{listings}
\lstdefinestyle{mystyle}{
    basicstyle=\footnotesize\ttfamily,
    breaklines=true,
    numbers=left,
    numbersep=5pt,
    showstringspaces=false
}
\lstset{style=mystyle}

% Поля страницы
\usepackage{geometry}
\geometry{
    left=2.5cm,
    right=1.5cm,
    top=2cm,
    bottom=2cm
}

%%%%%%%%%%%%%%%%%%%%%%%%%%%%%%%%%%%%%%%%%%%%%%%% Всякое для ГОСТа

% Для подражания выравниванию по ширине 
\usepackage[none]{hyphenat}  % полное отключение переносов
\usepackage{microtype}       % улучшает выравнивание (важно!)
\emergencystretch=5em   % разрешаем растягивать строки
\tolerance=500          % терпимость к плохим строкам

% Межстрочный интервал
\usepackage{setspace}
\onehalfspacing 

% Отступы в списках
\usepackage{enumitem}
\setlist{
    itemsep=0pt,      % расстояние между пунктами
    parsep=0pt,       % расстояние между абзацами внутри пункта
    topsep=0pt,       % расстояние до списка сверху
    partopsep=0pt,    % дополнительное расстояние
    leftmargin=*,     % отступ текста
    labelindent=\parindent,  % отступ маркера списка
    align=left         % для labelindent
}

%%%%%%%%%%%%%%%%%%%%%%%%%%%%%%%%%%%%%%%%%%%%%%%% НАСТРОЙКА ЗАГОЛОВКОВ ПО МЕТОДИЧКЕ
\usepackage{titlesec}

% 1. Глава (\section) всегда начинается с новой страницы
\newcommand{\sectionbreak}{\clearpage}

% 2. Формат главы: 16пт, жирный, ПРОПИСНЫЕ
\titleformat{\section}
  {\fontsize{16}{20}\selectfont\bfseries\raggedright\MakeUppercase} 
  {\thesection} 
  {1em} 
  {} 

% 3. Интервалы главы: ДО - 0pt (т.к. новая страница), ПОСЛЕ - 24pt
\titlespacing*{\section}
  {0pt}   % отступ слева
  {0pt}   % интервал ДО
  {24pt}  % интервал ПОСЛЕ

% 4. Формат подраздела (\subsection): 14пт, жирный
\titleformat{\subsection}
  {\fontsize{14}{18}\selectfont\bfseries\raggedright}
  {\thesubsection}
  {1em}
  {}

% 5. Интервалы подраздела: ДО - 24pt, ПОСЛЕ - 12pt
\titlespacing*{\subsection}
  {0pt}   % отступ слева
  {24pt}  % интервал ДО
  {12pt}  % интервал ПОСЛЕ
%%%%%%%%%%%%%%%%%%%%%%%%%%%%%%%%%%%%%%%%%%%%%%%%

% Переопределяем название листингов
\renewcommand{\lstlistingname}{Листинг}

% Меняем формат caption
\usepackage{caption}
\captionsetup[lstlisting]{labelsep=endash} % тире для листингов

% Как форматировать листинги (язык по умолчанию)
\lstset{
    language=C++,
    frame=single,
    captionpos=b,            % подпись снизу            
    numbers=left,            % нумерация слева (внутри рамки)
    numbersep=5pt,           % расстояние от номера до кода
    xleftmargin=15pt,        % отступ слева (вместе с рамкой)
    framexleftmargin=15pt,   % отступ рамки от левого края
    xrightmargin=5pt,
    framexrightmargin=5pt
}

% Формат подписей рисунков (точка на тире)
\DeclareCaptionFormat{myformat}{#1 -- #3}
\captionsetup[figure]{name=Рисунок,format=myformat}

% Формат подписей таблиц (точка на тире)
\captionsetup[table]{labelsep=endash} 

% Для ссылок по ГОСТу
\usepackage{cite} 

% Теоремы, определения и прочее
\newtheorem{theorem}{Теорема}[section]
\newtheorem{lemma}{Лемма}[section]
\newtheorem{definition}{Определение}[section]
\newtheorem{proposition}{Утверждение}[section]
\newtheorem{corollary}{Следствие}[section]

% Углубленная структура секций (до 4 уровня)
\newcommand{\subsubsubsection}[1]{\paragraph{#1}\mbox{}\\}
\setcounter{secnumdepth}{4}
\setcounter{tocdepth}{4}

% Гиперссылки
\usepackage{hyperref}
\usepackage{bookmark}

\usepackage{tikz}
\usetikzlibrary{shapes.geometric, arrows.meta, positioning, calc}

% Обзор компактный (3-4 листа): разделы идут слитно, а не с новой страницы.
\renewcommand{\sectionbreak}{\par\addvspace{2.5ex}}

\begin{document}

\begin{titlepage}
	\begin{center}
		\large
		{МИНИСТЕРСТВО НАУКИ И ВЫСШЕГО ОБРАЗОВАНИЯ\\ РОССИЙСКОЙ ФЕДЕРАЦИИ}

		Федеральное государственное автономное образовательное учреждение высшего образования
		\vspace{0.5cm}

		\textbf{<<Национальный исследовательский Нижегородский государственный университет им. Н.И. Лобачевского>>}\\
		(ННГУ)\\
		\vspace{1cm}
		\textbf{Институт информационных технологий, математики и механики}\\
		\vspace{0.5cm}
		\textbf{Кафедра: Ма­те­ма­ти­че­ско­го обес­пе­че­ния и су­пер­ком­пью­тер­ных тех­но­ло­гий}
		\vspace{1cm}

		Направление подготовки: <<Фундаментальная информатика и информационные технологии>>\\
		Профиль подготовки: <<Инженерия программного обеспечения>>
		\vfill

		\vfill

		\Large
		Аналитический обзор \\
		\normalsize
        по научно-исследовательской работе \\
		на тему:\\
		\Large
		\textbf{<<Методы улучшения поверхностных полигональных сеток биомедицинских моделей>>}
		{\LARGE
		}
		\bigskip


	\end{center}
	\vfill

	\hfill\begin{minipage}{0.4\textwidth}
		\textbf{Выполнил:}

            студент группы 3823Б1ФИ1

            \underline{\hspace{3cm}} Дорофеев И.\,Д. \bigskip

	\end{minipage}%

	\hfill\begin{minipage}{0.4\textwidth}
		\textbf{Научный руководитель:}

		Доцент кафедры ПМ, к.ф.-м.н.,

		\underline{\hspace{3cm}}  Кириллин М.\,Ю.

	\end{minipage}
	\vfill

	\begin{center}
		Нижний Новгород\\
		2026 г.
	\end{center}
\end{titlepage}
\thispagestyle{empty} % убираем колонтитулы с этой страницы

\newpage

% ==============================================
\section{Введение}

Поверхностные полигональные сетки, полученные извлечением изоповерхностей из сегментированных биомедицинских объёмов (МРТ, КТ, микроскопия), редко пригодны для дальнейшего использования в исходном виде. Они несут характерные дефекты: ступенчатые (terracing) артефакты вдоль оси стопки срезов, избыточную плотность (десятки миллионов граней), вырожденные и <<игольчатые>> треугольники, а также локальные не-манифолдные особенности и самопересечения. Для визуализации, конечно-элементного анализа и моделирования переноса излучения такие сетки требуют этапа \textbf{улучшения} --- согласованной обработки, повышающей качество элементов и корректность топологии при сохранении анатомической формы и объёма.

Настоящий обзор систематизирует методы улучшения поверхностных сеток применительно к биологическим данным по трём направлениям: сглаживание (устранение шума и ступенчатости), ремешинг и упрощение (нормализация связности и снижение сложности), а также появившиеся в последние годы обучаемые подходы. Сквозным требованием для биомедицины является \textbf{сохранение объёма и признаков}: избыточное сглаживание <<растворяет>> тонкие структуры, а агрессивное упрощение искажает границы тканей.

% ==============================================
\section{Сглаживание сеток}

Простейшее лапласово сглаживание сдвигает каждую вершину к центроиду соседей, но приводит к прогрессирующей усадке (shrinkage) --- модель <<тает>>. Метод Таубина \cite{taubin1995} трактует сглаживание как фильтрацию в спектре оператора Лапласа и чередует шаги с положительным и отрицательным коэффициентами ($\lambda\!\mid\!\mu$), подавляя высокочастотный шум без потери объёма. Неявное сглаживание \cite{desbrun1999} использует диффузию и поток по кривизне с безусловно устойчивым неявным шагом, что позволяет сильно сглаживать за одну итерацию. Улучшенное лапласово сглаживание HC \cite{vollmer1999} добавляет корректирующий член, возвращающий вершины к исходным позициям и компенсирующий усадку.

Ключевой для биомедицины является проблема \textbf{ступенчатых артефактов}: они порождаются осевой природой сканирования и не должны сглаживаться наравне с истинными деталями формы. Метод \texttt{Staircase-Aware Smoothing} \cite{munch2010} использует известное направление стопки срезов: по скалярному произведению нормалей граней с осью стопки классифицируются <<артефактные>> вершины (инцидентные одновременно параллельным и ортогональным оси граням), к которым применяется усиленное неоднородное сглаживание, тогда как истинные признаки сохраняются. Именно этот подход лежит в основе реализации, указанной в техническом задании. Родственные работы развивают адаптивное сглаживание с сохранением тонких признаков через кривизну и предотвращение самопересечений \cite{wang2017}.

% ==============================================
\section{Ремешинг и упрощение}

Второе направление нормализует связность и снижает сложность сетки. Упрощение методом стягивания рёбер на основе метрики квадратичной ошибки (QEM) \cite{garland1997} остаётся отраслевым стандартом: для каждой вершины накапливается матрица суммы квадратов расстояний до смежных плоскостей, и рёбра схлопываются в порядке возрастания ошибки, что сохраняет детализацию в областях высокой кривизны. Обзор методов упрощения полигональных моделей, в том числе с ускорением на графическом процессоре, дан в русскоязычной работе \cite{rus_gonahchyan}.

Изотропный ремешинг \cite{botsch2004} перестраивает сетку в почти равносторонние треугольники заданного размера атомарными операциями (расщепление, схлопывание и переброс рёбер, лапласова релаксация), что критично для устойчивости последующих расчётных алгоритмов. Полевые методы, такие как \texttt{Instant Field-Aligned Meshes} \cite{jakob2015}, выравнивают элементы по направлениям главных кривизн, порождая структурированные сетки. Отдельную задачу составляет \textbf{репарация}: пакет \texttt{MeshFix} \cite{attene2010} за счёт локальных операций гарантированно приводит сетку к водонепроницаемому манифолдному виду; современный метод репарации с сохранением признаков на основе ограниченной степенной диаграммы \cite{repair2025} строит манифолдные водонепроницаемые сетки, не разрушая характерные линии.

% ==============================================
\section{Обучаемые методы улучшения}

За последние годы для улучшения сеток активно привлекается машинное обучение. Сглаживание формулируется как задача обучения с подкреплением: модель на графовых нейросетях перемещает узлы в оптимальные положения, обучаясь на качественных сетках \cite{wang2023drl, gnnrl2024}. Лёгкие графовые модели \cite{gmsnet2024} и свёрточные сети \cite{cnn2025} достигают качества оптимизационного сглаживания при существенно меньших вычислительных затратах. Обучаемые подходы применяются и к упрощению: соответствующий обзор систематизирует методы на основе нейросетей \cite{simplifysurvey2025}. Исчерпывающий обзор геометрических и обучаемых методов подавления шума (denoising) на сетках приведён в \cite{denoising2023}. Наряду с обучаемыми развиваются и строгие оптимизационные методы: улучшение сеток с геометрическими ограничениями через коническое программирование второго порядка \cite{meshcone2024}.

% ==============================================
\section{Биомедицинская специфика и выводы}

Применительно к анатомическим моделям качество обработки напрямую влияет на достоверность последующего анализа. Исследование точности преобразования сетки в твердотельную модель методологией Тагути \cite{taguchi2023} показало, что наибольшее влияние на итоговую геометрию оказывает именно фактор сглаживания, тогда как снижение числа полигонов статистически менее значимо --- что подтверждает приоритет аккуратного сглаживания над агрессивной децимацией. Для сложных анатомических поверхностей эффективен ремешинг через сохраняющую площади параметризацию \cite{choi2021}, дающий однородные треугольники при контролируемом искажении.

Обобщая, надёжный конвейер улучшения биомедицинских сеток комбинирует: сглаживание с учётом ступенчатости и с сохранением объёма (Таубин, staircase-aware), изотропный ремешинг для нормализации связности, контролируемое QEM-упрощение и гарантированную репарацию водонепроницаемости. Обучаемые методы последних пяти лет перспективны как ускорители оптимизационного сглаживания, однако в задачах, требующих строгих топологических гарантий (конечно-элементный анализ, поверхностный перенос света), детерминированные методы с сохранением признаков и объёма остаются основой.

% ==============================================
\renewcommand{\refname}{Список литературы}
\addcontentsline{toc}{section}{Список литературы}
\begin{thebibliography}{99}

\bibitem{taubin1995} Taubin G. A signal processing approach to fair surface design // Proc. ACM SIGGRAPH. --- 1995. --- P. 351--358.

\bibitem{desbrun1999} Desbrun M., Meyer M., Schr\"oder P., Barr A. H. Implicit fairing of irregular meshes using diffusion and curvature flow // Proc. ACM SIGGRAPH. --- 1999. --- P. 317--324.

\bibitem{vollmer1999} Vollmer J., Mencl R., M\"uller H. Improved Laplacian smoothing of noisy surface meshes // Computer Graphics Forum (Eurographics). --- 1999. --- Vol. 18, no. 3. --- P. 131--138.

\bibitem{garland1997} Garland M., Heckbert P. S. Surface simplification using quadric error metrics // Proc. ACM SIGGRAPH. --- 1997. --- P. 209--216.

\bibitem{botsch2004} Botsch M., Kobbelt L. A remeshing approach to multiresolution modeling // Proc. Eurographics/ACM SIGGRAPH Symp. on Geometry Processing. --- 2004. --- P. 185--192.

\bibitem{jakob2015} Jakob W., Tarini M., Panozzo D., Sorkine-Hornung O. Instant field-aligned meshes // ACM Transactions on Graphics (SIGGRAPH Asia). --- 2015. --- Vol. 34, no. 6. --- Art. 189.

\bibitem{attene2010} Attene M. A lightweight approach to repairing digitized polygon meshes // The Visual Computer. --- 2010. --- Vol. 26, no. 11. --- P. 1393--1406.

\bibitem{munch2010} M\"unch T., Adler S., Preim B. Staircase-aware smoothing of medical surface meshes // Eurographics Workshop on Visual Computing for Biology and Medicine (VCBM). --- 2010. --- P. 83--90.

\bibitem{wang2017} Wang J., Yu Z. Quality improvement of surface triangular mesh using a modified Laplacian smoothing approach avoiding intersection // PLOS ONE. --- 2017. --- Vol. 12, no. 9. --- Art. e0184206.

\bibitem{choi2021} Choi G. P. T. et al. Adaptive area-preserving parameterization of open and closed anatomical surfaces // Computers in Biology and Medicine. --- 2022. (arXiv:2111.04265, 2021).

\bibitem{wang2023drl} Wang Z. et al. Unstructured surface mesh smoothing method based on deep reinforcement learning // Computational Mechanics. --- Springer, 2023. --- Vol. 72. --- P. 1249--1265.

\bibitem{gnnrl2024} GNNRL-smoothing: A prior-free reinforcement learning model for mesh smoothing // Journal of Computational Physics. --- Elsevier, 2025. (arXiv:2410.19834, 2024).

\bibitem{gmsnet2024} An intelligent mesh-smoothing method with graph neural networks // Frontiers of Information Technology \& Electronic Engineering. --- 2024.

\bibitem{cnn2025} A new mesh smoothing method based on convolutional neural network // Computational Mechanics. --- Springer, 2025.

\bibitem{denoising2023} Geometric and learning-based mesh denoising: A comprehensive survey // ACM Transactions on Multimedia Computing, Communications, and Applications. --- 2023. --- Art. 3625098.

\bibitem{repair2025} Feature-preserving mesh repair via restricted power diagram // Proc. ACM SIGGRAPH. --- 2025.

\bibitem{taguchi2023} Geometric accuracy analysis in mesh-to-solid conversion for 3D anatomical models through Taguchi methodology // Journal of Applied Research in Technology \& Engineering (JARTE). --- 2024.

\bibitem{simplifysurvey2025} A survey of learning-based mesh processing techniques for mesh simplification // Lecture Notes in Computer Science. --- Springer, 2025.

\bibitem{meshcone2024} MeshCone: Second-order cone programming for geometrically-constrained mesh enhancement // arXiv preprint arXiv:2412.08484. --- 2024.

\bibitem{rus_gonahchyan} Гонахчян В. И. Обзор методов упрощения полигональных моделей на графическом процессоре // Труды Института системного программирования РАН. --- 2014. --- Т. 26, вып. 2. --- С. 159--174.

\end{thebibliography}

\end{document}
